% !TeX program = xelatex
\documentclass[12pt,a4paper]{article}
\usepackage[margin=2cm]{geometry}
\usepackage{fontspec}
\usepackage{amsmath,amssymb}
\usepackage{enumitem}
\setmainfont{Times New Roman}
\setlength{\parindent}{0pt}
\begin{document}
\begin{center}
{\large\bfseries SỞ GIÁO DỤC VÀ ĐÀO TẠO HÀ NỘI}\\[4pt]
{\Large\bfseries ĐỀ THI VÀO LỚP 10 NĂM 2007}\\[2pt]
{\bfseries Môn: TOÁN -- PHẦN ĐẠI SỐ}\\
Thời gian làm bài: 90 phút
\end{center}
\textbf{Câu 1.} Cho biểu thức
\[
P=\frac{\sqrt{x}}{\sqrt{x}-1}+\frac{3}{\sqrt{x}+1}-\frac{6\sqrt{x}-4}{x-1}.
\]
\begin{enumerate}[label=\alph*)]
  \item Rút gọn $P$.
  \item Tìm các giá trị của $x$ để $P<\dfrac12$.
\end{enumerate}
\textbf{Câu 2.} Giải bài toán sau bằng cách lập phương trình:

Một người đi xe đạp từ $A$ đến $B$ cách nhau $24\,\mathrm{km}$. Khi từ $B$ trở về $A$, người đó tăng vận tốc thêm $4\,\mathrm{km/h}$ so với lúc đi, vì vậy thời gian về ít hơn thời gian đi 30 phút. Tính vận tốc của người đi xe đạp khi đi từ $A$ đến $B$.

\textbf{Câu 3.} Cho phương trình $x^2+bx+c=0$.
\begin{enumerate}[label=\alph*)]
  \item Giải phương trình khi $b=-3$, $c=2$.
  \item Tìm $b,c$ để phương trình có hai nghiệm phân biệt và tích hai nghiệm bằng $1$.
\end{enumerate}
\textbf{Câu 5.} Cho đường thẳng $y=(m-1)x+2$. Tìm $m$ để khoảng cách từ gốc tọa độ $O$ tới đường thẳng đó lớn nhất.
\end{document}
